\section{Introduction}
\label{sec:introduction}

A behavioral event log is a sparse record of when episodes of labeled activities
start and how long they last. The same log admits several readings at once. Read
at a one-minute resolution it is a binary indicator per label per slot, marking
the slots at which an episode begins. Read at ten minutes it is a count of
episodes together with a total occupied duration. Read at an hour or a day it is
a coarse summary that carries the trend and almost none of the detail. The
readings are not three datasets. They are one process observed through three
windows, and every quantity at a coarse window is a sum of quantities at a fine
one.

Three properties of such logs hold simultaneously and each of them is awkward for
a model defined at a single resolution. The log is sparse, because most slots at
the finest resolution contain nothing. It is organized at several granularities
at the same time, because the questions asked of it range from which minute an
episode began to whether a month differs from the month before. And the
dependencies among labels drift, because a routine that holds in one season is
not the routine of the next, so a graph fitted to the whole horizon describes an
average of several regimes and none of them.

Dynamic Bayesian networks are the natural starting point. They represent a
temporal stochastic process in factored form, admit latent variables, and keep
the temporal dependence interpretable \citep{dean1989model, murphy2002dbn}.
Learning their structure from data has been studied since structure search was
carried over to the dynamic case, hidden variables and incomplete data included
\citep{friedman1998learning}, and the score-based, constraint-based and hybrid
families developed for static networks all transfer to the lagged case
\citep{heckerman1995learning, spirtes2000causation, chickering2002optimal,
tsamardinos2006max}. What the standard formulation does not supply is a
commitment about how the resolutions relate to each other. A network fitted at
ten minutes and a network fitted at an hour are two models, and nothing in either
of them forces the first to sum to the second.

\subsection{Four requirements on a single model}
\label{sec:intro_requirements}

This work takes the position that the three readings should be one model, and
that such a model has to meet four requirements at the same time. Meeting three
of them is comparatively easy; the difficulty is that each of the four constrains
how the others can be met.

\paragraph{Exact aggregation between a fine count and its coarse summary.}
Let the resolutions be indexed by \(k \in \{1,\dots,K\}\), with \(k=1\) the
finest, let \(\gT^{(k)}\) be the ordered set of epochs at scale \(k\) and
\(T_k\) their number, and let the scales be nested through a parent map
\(\parmap_k\) with child sets \(\child_k(c) = \{t : \parmap_k(t) = c\}\). Writing
\(X_t^{(k,\ell)}\) for the count of label \(\ell \in \gL\) in epoch \(t\) at
scale \(k\) and \(Y_t^{(k,\ell)}\) for the total duration of that label in that
epoch, measured in finest-scale slots, the requirement is
\begin{align}
\sum_{t \in \child_k(c)} X_t^{(k,\ell)} = X_c^{(k+1,\ell)},
\qquad
\sum_{t \in \child_k(c)} Y_t^{(k,\ell)} = Y_c^{(k+1,\ell)},
\label{eq:intro_reconciliation}
\end{align}
for every \(k<K\), every parent epoch \(c \in \gT^{(k+1)}\) and every label. The
identity is a property of the data, since a coarse count is literally the sum of
the fine counts inside it. A model that predicts the two scales independently
will violate it, and repairing the violation after the fact is a separate
estimation problem. The n\(^2\)DBN removes the problem instead of solving it. The
coarsest scale is generated freely, and every finer scale is generated by
allocating the parent total among its children through a distribution whose
support is the set of non-negative integer vectors summing to that total, so
\cref{eq:intro_reconciliation} holds almost surely rather than approximately. In
the reported runs the measured reconciliation error is exactly zero for counts
and for durations in every fit of every method, including the fit on an observed
log.

\paragraph{Learned rather than assumed lagged dependencies.}
Which labels predict which, and at which lag, is the scientific content of the
graph, and fixing it by hand decides in advance what the analysis can find. The
model therefore carries a lagged parent set \(\gP_{k,\ell} \subseteq \gL \times
\{1,\dots,p_k\}\) per label and per scale, and learns it by maximizing a local
score penalized by \(\lambda_k |\gP_{k,\ell}|\) inside a bounded candidate set.
Because every admissible parent is lagged, any graph over the unrolled variables
is acyclic, so the search needs no acyclicity constraint and the score decomposes
over targets.

\paragraph{Separate representation of persistent drift and same-epoch
co-movement.}
Two labels can move together for two quite different reasons. One is a persistent
change in the process itself, a regime that raises or lowers the rate of several
labels for as long as it lasts. The other is a contemporaneous coupling within a
single epoch that carries no direction. Neither is a lagged dependence, and a
model that offers only lagged edges will report both of them as edges. The
n\(^2\)DBN gives each its own latent variable. A regime \(S_t^{(k)} \in
\{1,\dots,M_k\}\) follows a context-dependent Markov chain and carries persistent
drift, and a factor \(H_t^{(k)} \in \sR^{r_k}\) carries same-epoch co-movement.
No observed variable is a parent of another observed variable in the same epoch.
Contemporaneous association is mediated entirely through \(S_t^{(k)}\) and
\(H_t^{(k)}\), which is what keeps co-movement out of the reported graph and, at
the same time, keeps the unrolled graph acyclic.

\paragraph{A mapping back to an event trace with wall-clock times.}
A multiscale summary that cannot be turned back into events answers a narrower
question than the one that was asked. At the finest scale the model constrains
counts to be binary, so \(X_t^{(1,\ell)} = 1\) means an episode of label
\(\ell\) begins in slot \(t\), and \(Y_t^{(1,\ell)}\) is that episode's length in
slots. With a slot width \(\Delta_1\) and a time origin \(t_0\), the episode
beginning at slot \(t\) occupies the wall-clock interval
\begin{align}
\left[\,
t_0 + (t-1)\,\Delta_1,\;
t_0 + \bigl(t - 1 + Y_t^{(1,\ell)}\bigr)\,\Delta_1
\,\right),
\label{eq:intro_episode}
\end{align}
and the trace of label \(\ell\) is the sequence of these intervals in increasing
order of start slot. The map is deterministic, so a coarse-scale prediction
propagates to a dated event list without a second estimation step.

\subsection{Contributions}
\label{sec:intro_contributions}

The paper makes five contributions.

First, it formulates the nested non-stationary dynamic Bayesian network, a single
probabilistic model of event counts and durations at nested resolutions in which
\cref{eq:intro_reconciliation} holds by construction, lagged structure is learned
per scale, persistent drift and same-epoch co-movement are carried by separate
latent variables, and \cref{eq:intro_episode} returns a dated event trace.

Second, it derives a bidirectional variational Structural EM procedure for the
model, in which the regime posterior is obtained by forward-backward recursion,
the factor posterior by a Gaussian approximation, the parameters by monotone
block updates, and the graph by a penalized local search whose score is an
expectation under the current posterior over the latent variables and includes
the block normalizer of the allocation. One sweep is two passes, fine to coarse
and then coarse to fine, and each pass identifies the component of a scale-$k$
predictor the other leaves unidentified.

Third, it measures what that latent-conditional score buys, by placing the
structure step behind a common interface with published structure learners and
running them on the same generated data, inside the same in-degree budget, with
the same parameters fitted afterwards, so that the only difference between two
rows of a results table is the graph. The benchmark is screened beforehand, so
that a low score is attributable to an estimator rather than to a scenario
containing nothing to find. The screen also fixes how many true parents each
scale carries. It places lagged edges at the coarsest and at the middle scale
and none at the finest, where a lagged edge on a binary start process can only
move which slot of a block begins an episode and does not pay for itself, so
the finest scale grades a method on staying silent.

Fourth, it checks the same code against external ground truth, using released
Bayesian networks tiled into two-slice temporal networks, and establishes what
such a check requires. A latent-variable model carries free quantities that a
latent-free learner does not, namely the number of regimes, the factor
dimension and the edge penalty. Fixing them by hand and then reporting a
comparison measures the choices, not the criteria. We therefore select the
regime count and the edge penalty on held-out log evidence from the training
split alone, holding the factor dimension at zero because the tiled networks
carry no same-slice factor, and we show that on those networks the selection
recovers the setting that maximizes $F_1$ without ever reading the true graph.
On the generated grid the same rule picks a different edge penalty from the one
that maximizes $F_1$, and rescoring the frontier at every scale that carries a
true edge shows this is a property of the selector rather than of the metric:
the middle scale, the coarsest scale and the mean over scales each disagree with
held-out evidence, and the two truth-bearing scales disagree with each other, so
no single penalty is best at both.

Fifth, it documents the failures. A residual leak of same-epoch structure into
finest-scale lagged self edges survives every mechanism excluded for it; the
mean-field evaluation of the allocation normalizer is contained by an outer
guard rather than removed; recall on external data remains below that of hill
climbing after selection; and parameter recovery does not meet the sample-size
criterion stated before the experiment was run, from a ground-truth start as
well as from a data-driven one.
